\bookchapter{Bab 41: Panorama Alat Pengaturcaraan AI 2026}

\chapterintro{Peringatan dahulu: bab ini ialah satu “imej sepintas” ekosistem alat pada Ogos 2026. Enam bulan kemudian ia akan lapuk — dan itulah pelajaran pertama yang sebenarnya ingin diajarkan bab ini. Tujuan bab ini bukan mengajar anda “gunakan alat yang mana”, tetapi membantu anda membina kerangka pertimbangan “cara memilih alat”. Apabila anda memahami falsafah reka bentuk alat yang berbeza, anda boleh membuat pilihan yang sesuai untuk diri anda pada bila-bila masa.}

\booksection{Tiga Aliran Besar dan Kerangka Pemilihan}

Alat pengaturcaraan AI tahun 2026 secara amnya terbahagi kepada tiga aliran besar, dan di sebalik setiap aliran terdapat satu falsafah reka bentuk yang berbeza. Lihat dahulu gambaran keseluruhan ini:

\begin{center}
\begin{tabularx}{\linewidth}{|X|X|X|X|}
\hline
\textbf{Aliran} & \textbf{Alat utama} & \textbf{Falsafah reka bentuk} & \textbf{Senario sesuai} \tabularnewline
\hline
AI Terbina dalam IDE & Cursor 3, Trae, Windsurf, CodeBuddy & Mengintegrasikan keupayaan AI secara mendalam ke dalam penyunting & Pembangunan harian, semakan kod, penambahbaikan berperingkat \tabularnewline
Agent Terminal & Claude Code, Codex CLI, Aider & Terminal sebagai antara muka; AI mengendalikan sistem fail dan arahan secara terus & Operasi tahap projek, pengubahsuaian pukal, integrasi CI/CD \tabularnewline
Builder Awan & v0, Bolt, Replit Agent & Membina projek lengkap dari sifar dalam pelayar & Prototaip pantas, mula tanpa konfigurasi, pratonton serta-merta \tabularnewline
\hline
\end{tabularx}
\end{center}

\begin{center}\small\color{gray}Tiga aliran besar alat pengaturcaraan AI\end{center}

Kunci memilih alat bukan “yang mana terbaik”, tetapi “yang mana paling sesuai dengan senario semasa anda”. Tiga dimensi pertimbangan yang praktikal:

\begin{itemize}
  \item Bagaimana asas pengaturcaraan anda? Tanpa asas pilih Builder Awan, ada asas pilih AI Terbina dalam IDE, pengguna mahir pilih Agent Terminal.
  \item Di fasa manakah projek anda? Mula dari sifar pilih Builder Awan, pembangunan berulang pilih AI Terbina dalam IDE, penyelenggaraan dan penstrukturan semula pilih Agent Terminal.
  \item Apakah keperluan pelancaran anda? Pengesahan pantas pilih Builder Awan, projek serius pilih AI Terbina dalam IDE, aliran kerja automatik pilih Agent Terminal.
\end{itemize}

\booksection{Aliran AI Terbina dalam IDE}

AI Terbina dalam IDE ialah aliran paling arus perdana: keupayaan AI dipasang ke dalam penyunting kod anda, menambah baik dengan AI sambil mengekalkan pengalaman pengaturcaraan tradisional. Pada 2026, perwakilan aliran ini ialah Cursor 3, Trae, Windsurf dan CodeBuddy.

\booksubsection{Cursor 3}

Cursor kini ialah IDE pengaturcaraan AI yang paling luas digunakan di dunia. Cursor 3 yang dikeluarkan pada April 2026 menandakan satu titik peralihan penting bagi alat pengaturcaraan AI: kadar penggunaan mod Agent buat pertama kalinya melebihi mod pelengkapan Tab — pembangun beralih daripada “AI melengkapkan kod saya” kepada “AI menyiapkan tugasan saya”. Ciri teras Cursor 3 termasuk: penyuntingan pelbagai fail (AI boleh mengubah suai beberapa fail serentak untuk menyiapkan tugasan), kesedaran konteks (memahami struktur projek dan hubungan kebergantungan secara automatik) serta integrasi mendalam dengan terminal (AI boleh menjalankan arahan dan melihat keluaran terus dalam penyunting).

\booksubsection{Trae}

Trae ialah IDE pengaturcaraan AI yang dikeluarkan oleh ByteDance dan mempunyai pangkalan pengguna besar di pasaran China. Mod SOLO-nya (juga dipanggil mod Builder) ialah keupayaan khas “bina projek dengan satu ayat” — anda menerangkan keperluan, dan AI menyiapkan keseluruhan aliran kerja daripada pemulaan projek hingga penjanaan kod secara autonomi. Kelebihan Trae terletak pada sokongan bahasa Cina dan penyesuaian dengan persekitaran rangkaian di China, sesuai untuk pengguna di China bermula dengan pantas.

\booksubsection{Alat Lain}

Windsurf (dahulunya Codeium) terkenal dengan pengalaman pelengkapan kod yang lancar, sesuai untuk pembangun yang inginkan pengalaman “menaip secepat terbang”. Tencent Cloud CodeBuddy, sebagai sambungan VS Code, mahir berintegrasi dengan ekosistem awan dan sesuai untuk projek yang berancang melancarkan ke Tencent Cloud atau membangunkan mini program WeChat. Alibaba Tongyi Lingma dan Baidu Comate turut menawarkan pengalaman integrasi IDE yang serupa, masing-masing dengan keistimewaan dalam pengoptimuman senario bahasa Cina.

\booksection{Aliran Agent Terminal}

Agent Terminal ialah aliran yang berkembang paling pesat pada 2026. Idea terasnya: AI tidak sepatutnya sekadar sambungan penyunting; ia sepatutnya dapat mengendalikan keseluruhan persekitaran pembangunan seperti seorang pembangun — membaca dan menulis fail, menjalankan arahan, melihat keluaran, memasang kebergantungan. Claude Code ialah perwakilan aliran ini.

\booksubsection{Claude Code}

Claude Code ialah alat pengaturcaraan AI terminal yang dikeluarkan oleh Anthropic, dan menjelang pertengahan 2026 telah mencapai pendapatan tahunan AS\$2.5 bilion. Cara kerjanya berbeza sama sekali daripada IDE tradisional: anda menerangkan tugasan dalam bahasa semula jadi di terminal, dan Claude Code akan membaca fail berkaitan, mengubah suai kod, menjalankan ujian, melihat keputusan dan memperbaiki secara berulang — keseluruhan proses tidak memerlukan anda membuka penyunting.

Kelebihan Claude Code terletak pada “kesedaran keseluruhan projek”: ia tidak terhad oleh tetingkap penyunting, boleh mengimbas seluruh direktori projek dengan bebas dan memahami kebergantungan antara fail. Ia juga menyokong fail memori projek CLAUDE.md, supaya AI memahami latar belakang projek pada setiap permulaan sesi. Bagi tugasan seperti penstrukturan semula projek besar, pengubahsuaian pukal dan penulisan ujian, Agent Terminal lebih cekap daripada sambungan IDE.

\booksubsection{Codex CLI dan Aider}

Codex CLI daripada OpenAI menawarkan pengalaman pengaturcaraan terminal yang serupa dan menyokong fungsi /import — mengimport kod sedia ada untuk difahami oleh AI. Aider ialah alat pengaturcaraan AI terminal sumber terbuka yang menyokong pelbagai model (GPT, Claude, DeepSeek dan lain-lain), sesuai untuk pengguna yang mahu bertukar model secara fleksibel. Walaupun setiap alat ini ada keistimewaannya sendiri, idea terasnya sama: membolehkan AI menguasai sepenuhnya persekitaran pembangunan dalam terminal.

\booksection{Aliran Builder Awan}

Builder Awan ialah aliran paling sesuai untuk pengguna tanpa asas. Anda tidak perlu memasang sebarang perisian; cukup buka pelayar, terangkan keperluan anda, dan AI akan membina projek lengkap di awan serta menyediakan pratonton serta-merta.

\begin{itemize}
  \item v0 (dibina oleh Vercel): mahir menjana komponen dan halaman UI bahagian hadapan, sesuai untuk membuat prototaip laman web dengan pantas.
  \item Bolt: menyokong penciptaan aplikasi penuh (full stack) daripada satu ayat, merangkumi bahagian hadapan, bahagian belakang dan pangkalan data.
  \item Replit Agent: bekerja dalam IDE awan Replit; bukan sahaja boleh membina projek, malah boleh menjalankan dan melancarkan terus di awan.
\end{itemize}

Kelebihan bersama alat-alat ini ialah “mula tanpa konfigurasi” — cukup buka pelayar untuk bermula. Keterbatasannya: apabila projek menjadi rumit, fleksibiliti persekitaran penyuntingan awan tidak sebaik IDE tempatan. Oleh itu, ia paling sesuai untuk fasa prototaip “dari 0 ke 1”: gunakan Builder Awan untuk mengesahkan idea dengan pantas, dan selepas pengesahan berjaya, pindahkan ke IDE tempatan untuk pembangunan mendalam.

\booksection{Alat Berubah, Metodologi Tidak Berubah}

Andai esok semua alat pengekodan AI hilang, tiga tabiat yang tetap perlu anda kekalkan: fikir jernih dahulu sebelum bertindak, buat satu langkah kecil pada satu-satu masa, dan sahkan selepas siap. Alat ialah pembesar, metodologi ialah asas. Pembesar boleh ditukar, asas tidak boleh digugurkan.

Menyemak perubahan alat daripada 2024 hingga 2026: pada 2024 GitHub Copilot ialah raja, pada 2025 Cursor mengambil alih, dan pada 2026 Claude Code bangkit mengejut. Setiap kali alat baharu muncul, ada yang mendakwa “kali ini berbeza”, tetapi akhirnya didapati: mereka yang mahir menggunakan alat baharu justru mereka yang mempunyai metodologi kukuh. Mereka hanya memerlukan beberapa hari untuk menyesuaikan diri apabila bertukar alat, manakala mereka yang hanya belajar mengendalikan alat terpaksa belajar semula setiap kali bertukar.

Oleh itu, laluan pembelajaran yang disyorkan buku ini ialah: mula-mula belajar dengan satu alat (alat mana pun boleh) dan rasai lima prinsip pemikiran utama melalui amalan; apabila metodologi anda cukup kukuh, barulah cuba alat daripada aliran yang berbeza dan rasai perbezaan falsafah reka bentuknya; akhirnya, anda akan membentuk gabungan alat anda sendiri — contohnya, guna Cursor untuk pembangunan harian, Claude Code untuk penstrukturan semula dan ujian, serta v0 untuk prototaip pantas. Gabungan alat berbeza bagi setiap orang, tetapi metodologinya sama.

\begin{tcbpractice}{Mari praktik: Sediakan satu set "senjata gabungan" untuk diri anda}
Gunakan tiga dimensi pertimbangan (asas pengaturcaraan, fasa projek, keperluan pelancaran) untuk menentukan kedudukan anda: jika anda tanpa asas cuba Bolt atau v0 untuk membuat satu halaman; ada asas guna Cursor atau Trae untuk satu tugasan berulang; mahu mencabar diri jalankan Claude Code sekali di terminal. Kemudian jawab: aliran manakah paling sesuai dengan fasa anda sekarang?

\end{tcbpractice}

\begin{tcbsummary}{Ringkasan Bab Ini}
\begin{itemize}
  \item Tiga aliran besar: AI Terbina dalam IDE (pembangunan harian), Agent Terminal (operasi tahap projek), Builder Awan (prototaip untuk pemula)
  \item Memilih alat dengan tiga soalan: asas pengaturcaraan? fasa projek? keperluan pelancaran?
  \item Mod Agent Cursor 3 buat pertama kalinya mengatasi pelengkapan Tab; Claude Code mencapai pendapatan tahunan AS\$2.5 bilion
  \item Builder Awan sesuai untuk "dari 0 ke 1"; selepas pengesahan, pindahkan semula ke IDE tempatan untuk pembangunan mendalam
  \item Alat bertambah baik setiap tiga bulan: mereka yang menguasai metodologi hanya memerlukan beberapa hari penyesuaian apabila bertukar alat
\end{itemize}

\end{tcbsummary}
